\section{Resource management}
\label{sec:resource}

This section names who owns each package, puts those people on the Table~\ref{tab:precedence} dates, and shows where demand would pass twelve and what we do about it.

\subsection{Roles and responsibilities}
\label{sec:resource-raci}

Table~\ref{tab:raci} is the RACI for this campaign. Responsible does the work. Accountable answers for it. Consulted sees the work before the Accountable owner signs. Each Level-3 package has exactly one Accountable owner.

There are two layers, both specific to Bowen.

\textbf{Consultancy layer:}
\begin{itemize}[leftmargin=1.4em,itemsep=0.25em]
\item Abdul Wahab Obeid (project manager and integration).
\item Hadi Alkhub (planner).
\item Ziyad Alahmari (risk, quality and HSE).
\item Ilsa Hashmi (cost, resources and the Juru commercial interface).
\end{itemize}

\textbf{Field layer:}
\begin{itemize}[leftmargin=1.4em,itemsep=0.25em]
\item Launch Director.
\item Range Safety Officer.
\item GSE Lead.
\item Avionics Specialist.
\item Juru Cultural Officer.
\item Design-and-construct package manager.
\end{itemize}

The owners follow the network.
\begin{itemize}[leftmargin=1.4em,itemsep=0.25em]
\item Hadi is Accountable for stacking because the dawn-lift window is a schedule control (R-06).
\item The Launch Director is Accountable for static fire and lift-off, with the RSO consulted.
\item Range exclusion is the other way round: the RSO is Accountable and the Launch Director is Responsible.
\item Closeout audit is Ziyad's, so quality is not marking its own homework.
\end{itemize}

\begin{table}[H]
\centering
\footnotesize
\setlength{\tabcolsep}{3.6pt}
\caption[RACI for Level-3 packages]{RACI for Level-3 packages. Exactly one Accountable owner per row. Field titles are as used on the pad; consultancy names are the Group~1 officers.}
\label{tab:raci}
\begin{tabularx}{\textwidth}{@{}l >{\raggedright\arraybackslash}p{3.15cm} >{\raggedright\arraybackslash}p{2.55cm} >{\raggedright\arraybackslash}p{2.35cm} X@{}}
\toprule
WBS & Package & A & R & C \\
\midrule
1.1.1 & PEP and integration & Abdul & Hadi & Ziyad, Ilsa \\
1.1.2 & Governance and EVM & Ilsa & Hadi & Abdul \\
1.1.3 & ASA / CASA / GBRMPA & Abdul & Hadi & Launch Director \\
1.1.4 & Juru heritage & Juru Officer & Ilsa & Abdul, Ziyad \\
1.2.1 & Receiving / cleanroom & GSE Lead & D\&C mgr & Hadi, Ilsa \\
1.2.2 & Crane and stacking & Hadi & GSE Lead & Launch Director, RSO \\
1.2.3 & Cryo GSE & D\&C mgr & GSE Lead & Ziyad \\
1.2.4 & RF calibration & Avionics & Hadi & Launch Director \\
1.3.1 & Avionics power-on & Avionics & GSE Lead & Hadi \\
1.3.2 & WDR (cold) & GSE Lead & Juru Officer & Ziyad, LD, RSO \\
1.3.3 & Static fire & Launch Director & GSE Lead & RSO, Ziyad, Avionics \\
1.3.4 & Post-fire / pre-kitted spare & Avionics & GSE Lead & Ziyad \\
1.4.1 & LRR and RSO & Launch Director & RSO & Abdul, Ziyad, Juru Officer, payload sponsor \\
1.4.2 & Range exclusion & RSO & Launch Director & Juru Officer \\
1.4.3 & Launch ($T$-0) & Launch Director & RSO, Avionics & Abdul \\
1.4.4 & Telemetry handover & Avionics & Hadi & Ilsa \\
1.4.5 & Demob and audit & Ziyad & D\&C mgr & Ilsa, Abdul \\
\bottomrule
\end{tabularx}
\end{table}

\subsection{Resource allocation and the schedule}
\label{sec:resource-capacity}

Every bar in Table~\ref{tab:precedence} carries a named pad crew. The standing crew is eight: four GSE technicians, two stacking technicians, the Launch Director and the RSO. Juru monitors, avionics and RF specialists join when their bars are live. Eight is the standing organisation, not a hard pad limit.

The hard limit is twelve people on the pad at once, from the blast zone, suit air and access control. Figure~\ref{fig:resources} must not cross it. These occupancy numbers are PEP controls, not Week~4 charter figures. Demand comes from the schedule: each day, add the people on every live bar and count each person once. The weekly bar is that week's busiest day. Twelve is a same-time limit, not a weekly headcount.

\subsection{Where demand peaks, and what we do}

October is the peak because stacking, Wet Dress Rehearsal and static fire sit there. Leaving A-125 on its late dates (7--24~October) and adding four overtime heads on 16--30~October asks for fourteen people on 16--24~October: two over the cap of twelve, and six over the crew of eight. Week~9 is already thirteen, because late RF is still on the pad during stacking. After 24~October the RF specialists are gone, so leftover overtime sits on static fire. Figure~\ref{fig:resources} and Table~\ref{tab:resource-build} show that over-allocation. It is R-13.

Four responses were tested: smoothing, levelling, hiring and resequencing. The baseline is smoothing, then hiring. Levelling and resequencing are rejected.

Smoothing moves work only inside its float, so the finish date stays. A-125 has 39 days of total float and is not launch-critical. Placed on its early dates, 29~August to 15~September, the two RF specialists sit in weeks 3--6. Cost is \$0. Lift-off stays 15~November. Static fire on 25--30~October still needs more than eight standing heads. Wet Dress Rehearsal on 16--24~October is inside the cap once RF is gone, and it does not share a day with stacking.

Levelling would delay $TF = 0$ work. Wet Dress Rehearsal and static fire are already $TF = 0$, so pushing either pushes 15~November. Levelling is rejected.

Resequencing would start the rehearsal before stacking finishes, or static fire before permits and RF are done, and would break the hold points in Section~\ref{sec:schedule}. RF cannot follow static fire: the array has to be live before hot fire. Resequencing is rejected.

Hiring adds four surge technicians for seven days, 25--31~October, inside WBS~1.3.3 at \$100/h: $4 \times 7 \times 8 \times \$100 = \$22{,}400$, already in the \$1{,}500{,}000 estimate \citep{fwc2026}. They replace the overtime heads. Mitigated 25--30~October is eight standing heads plus four surge: twelve, on the cap.

A ten-hour stagger on Wet Dress Rehearsal and static fire is the third R-13 control. Named people can share the calendar day without sharing the blast zone in the same hour. Cost is \$0 and no date moves. Figure~\ref{fig:resources} counts unique people that day, not the hour, so the figure does not change. Option~C, if invoked, shortens post-fire diagnostics only and does not redraw it.

\begin{figure}[H]
\centering
\includegraphics[width=0.98\linewidth]{resource_histogram.pdf}
\caption[Weekly pad demand versus capacity]{Weekly pad demand versus capacity, from Table~\ref{tab:resource-build}. Each bar is the busiest day that week. Hashed crimson is the unmitigated series (peak 14 on 16--24~October): Wet Dress Rehearsal, RF still on its late bar, and four overtime heads. Teal is the mitigated baseline (peak 12): RF moved to weeks 3--6, and four surge technicians on 25--31~October only. The navy dashed line is the standing crew of eight. The crimson dotted line is the twelve-person pad cap. Eight is not a hard pad limit.}
\label{fig:resources}
\end{figure}

% Generated by scripts/plot_resource_figures.py — do not edit by hand.
\begin{table}[H]
\centering
\footnotesize
\setlength{\tabcolsep}{3.2pt}
\caption[Weekly simultaneous pad occupancy]{Weekly simultaneous pad occupancy, from activity $\times$ named crew. Each weekly total is the maximum unique-head count on any calendar day that week (HSE cap is simultaneous, not weekly unique). Mitigated A-125 uses $ES$--$EF$; unmitigated uses the late bar $LS$--$LF$. Surge is four heads on 24--30~Oct only (7~days). Unmitigated overtime is four heads on WDR and static fire (16--30~Oct). Base, RF, surge and OT are disjoint name sets, so they add on the peak day.}
\label{tab:resource-build}
\begin{tabular}{@{}r l r r r r l r r r r@{}}
\toprule
Week & Mit.\ day & Base & RF & Surge & Mit. & Unmit.\ day & Base & RF & OT & Unmit. \\
\midrule
1 & 10~Aug & 6 & 0 & 0 & 6 & 10~Aug & 6 & 0 & 0 & 6 \\
2 & 17~Aug & 6 & 0 & 0 & 6 & 17~Aug & 6 & 0 & 0 & 6 \\
3 & 29~Aug & 6 & 2 & 0 & 8 & 24~Aug & 6 & 0 & 0 & 6 \\
4 & 31~Aug & 6 & 2 & 0 & 8 & 31~Aug & 6 & 0 & 0 & 6 \\
5 & 7~Sep & 6 & 2 & 0 & 8 & 7~Sep & 6 & 0 & 0 & 6 \\
6 & 14~Sep & 8 & 2 & 0 & 10 & 14~Sep & 8 & 0 & 0 & 8 \\
7 & 22~Sep & 11 & 0 & 0 & 11 & 22~Sep & 11 & 0 & 0 & 11 \\
8 & 28~Sep & 11 & 0 & 0 & 11 & 28~Sep & 11 & 0 & 0 & 11 \\
9 & 5~Oct & 11 & 0 & 0 & 11 & 7~Oct & 11 & 2 & 0 & 13 \\
10 & 12~Oct & 9 & 0 & 0 & 9 & 16~Oct & 8 & 2 & 4 & 14 \\
11 & 24~Oct & 8 & 0 & 4 & 12 & 19~Oct & 8 & 2 & 4 & 14 \\
12 & 26~Oct & 8 & 0 & 4 & 12 & 26~Oct & 8 & 0 & 4 & 12 \\
13 & 6~Nov & 5 & 0 & 0 & 5 & 6~Nov & 5 & 0 & 0 & 5 \\
14 & 15~Nov & 8 & 0 & 0 & 8 & 15~Nov & 8 & 0 & 0 & 8 \\
15 & 21~Nov & 4 & 0 & 0 & 4 & 21~Nov & 4 & 0 & 0 & 4 \\
16 & 23~Nov & 4 & 0 & 0 & 4 & 23~Nov & 4 & 0 & 0 & 4 \\
17 & 30~Nov & 4 & 0 & 0 & 4 & 30~Nov & 4 & 0 & 0 & 4 \\
\bottomrule
\end{tabular}
\end{table}


\subsection{Resource-conflict escalation}
\label{sec:resource-escalation}

The live conflict is GSE technicians wanted on Wet Dress Rehearsal and on static fire in late October. The path is timed and named. It is not a values statement.

\begin{itemize}[leftmargin=1.4em,itemsep=0.25em]
\item Tier~1 is the field lead. They have four hours for a same-day fix: the ten-hour stagger, or borrow from float, with no extra spend. That is not calendar levelling. It does not delay a $TF = 0$ bar and it does not move 15~November.
\item Tier~2 is the campaign PM (Abdul). They have twelve hours to choose overtime, hire or resequence if the stagger is not enough. Resequencing here is a live trade-off, not the baseline: the baseline already rejected changing the order of stacking, permits and hot fire. The cost goes to Ilsa so that Table~\ref{tab:wbs-cost} stays the book of record.
\item Tier~3 is the Executive Board. They have 24 hours. They sit only if the choice would spend contingency, or would move static fire (M-5, 30~October) or lift-off (M-7, 15~November). Those dates are the same dates as Section~\ref{sec:schedule}.
\end{itemize}
